\chapter{Conclusão}
\label{cap:conclusao-redes}

Este manual percorreu a rede do Linux sob a ótica do eBPF. Os capítulos
\ref{cap:fundamentos} a \ref{cap:tc} estabeleceram a base: onde a pilha de rede
oferece ganchos, o que o XDP e o TC podem ver e fazer em cada um deles, e como mapas
BPF levam informação e configuração entre o kernel e o espaço de usuário. O capítulo
\ref{cap:obs-rede} usou essa base para observar tráfego sem alterá-lo, e o capítulo
\ref{cap:seguranca-rede} a usou para agir sobre ele, com listas de bloqueio em mapas e
a limitação de SYN por origem. A progressão se resume em uma regra prática que se
repete ao longo do texto: observar primeiro, impor depois.

Os capítulos \ref{cap:k8s} e \ref{cap:cilium} levaram as mesmas ideias a contêineres e
ao Kubernetes, onde a identidade de uma carga de trabalho substitui o endereço IP
como critério de política, e o Cilium e o Hubble fazem em escala de cluster o que os
exemplos de laboratório fazem em uma interface. O capítulo \ref{cap:ferramentas-rede}
mostrou o conjunto de ferramentas (BCC, \textit{bpftrace}, \textit{bpftool} e
\textit{libbpf}) e o capítulo \ref{cap:exemplos-redes} as reuniu em três cenários
executáveis. Vale reforçar a distinção feita em cada capítulo: o que é conceitual
(Kubernetes e Cilium, sem cluster disponível) e o que é executável (o restante), e, entre
os executáveis, o que já foi compilado e o que aguarda execução com privilégio de
administrador.

O caminho natural para fechar essa lacuna está na pasta \texttt{testes/} do repositório.
O script \texttt{09-redes-exemplos.sh} executa os exemplos do capítulo
\ref{cap:exemplos-redes} em um laboratório com \textit{network namespace}, e a saída
guardada em \texttt{testes/resultados/} é a evidência que deve substituir as marcas
\texttt{\% TODO: validar com sudo} do texto. Para o leitor que deseja avançar, os
próximos passos são validar os capítulos de Kubernetes e Cilium em um cluster local de
teste, ampliar o \textit{firewall} para IPv6 e rastreamento de conexões, e conectar a
telemetria de rede à cadeia de Prometheus e Grafana apresentada no Manual~1.
